\chapterbackmatter{Supplementary Exercises}

\begin{enumerate}
\SupplementaryExercise{1}{The Klein--Gordon Current Is Not a Probability}
{(Adapted from H. Liu, MIT 8.323, Assignment 2, 2023,
Problem 1, pp. 2--3.)}{mit-823-s23-ps2}
For a complex solution of
\((\Box-m^2)\phi=0\), derive a conserved current and evaluate its density
on positive- and negative-frequency plane waves.  Explain why it is a charge
current rather than a probability current.

\SupplementaryExercise{2}{Recovering Scalar Mode Operators}
{(Adapted from H. Liu, MIT 8.323, Assignment 2, 2023,
Problem 2(b)--(c), pp. 3--4.)}{mit-823-s23-ps2}
At \(t=0\), invert the real scalar mode expansion to express
\(a(\mathbf p)\) and \(a^\dagger(\mathbf p)\) in terms of the spatial
Fourier transforms of \(\phi\) and its conjugate momentum.  Derive the
oscillator commutator from the canonical equal-time commutators.

\SupplementaryExercise{3}{Four-Momentum of a Free Scalar Field}
{(Adapted from H. Liu, MIT 8.323, Assignment 2, 2023,
Problem 3(a)--(b), p. 4.)}{mit-823-s23-ps2}
Starting from the symmetric stress tensor of a real scalar field, show after
normal ordering that
\[
 P^\mu=\int\frac{d^3p}{(2\pi)^3}\,
 p^\mu a^\dagger(\mathbf p)a(\mathbf p).
\]
Verify its action on a one-particle state.

\SupplementaryExercise{4}{Translations Generated by \(P^\mu\)}
{(Adapted from H. Liu, MIT 8.323, Assignment 2, 2023,
Problem 3(c), p. 4.)}{mit-823-s23-ps2}
Use the mode expression for \(P^\mu\) to prove
\([P^\mu,\phi(x)]=-i\partial^\mu\phi(x)\), and exponentiate the
commutator to obtain the finite translation law.

\SupplementaryExercise{5}{Orbital Lorentz Charge and Spin Zero}
{(Adapted from H. Liu, MIT 8.323, Assignment 2, 2023,
Problem 4, pp. 4--5.)}{mit-823-s23-ps2}
Derive the Lorentz current
\(M^{\rho\mu\nu}=x^\mu T^{\rho\nu}-x^\nu T^{\rho\mu}\)
for a free real scalar.  Put its rotation generators in momentum space and
show that there is no intrinsic spin term.

\SupplementaryExercise{6}{The Vacuum--Particle Field Matrix Element}
{(Adapted from H. Liu, MIT 8.323, Assignment 2, 2023,
Problems 2--3, pp. 3--4.)}{mit-823-s23-ps2}
With relativistically normalized states
\(\ket p=\sqrt{2p^0}\,a^\dagger(\mathbf p)\ket\Omega\), compute
\(\bra\Omega\phi(x)\ket p\).  Check its Klein--Gordon equation and
Lorentz transformation law.

\SupplementaryExercise{7}{The Lorentz Algebra in the Vector Representation}
{(Adapted from B. Mistlberger and K. Zhou, Stanford PHYS-330,
Problem Set 2, 2022, Problem 1, p. 1.)}{zhou-stanford-phys330-ps2-2022}
Construct the \(4\times4\) generators \(J_i\) and \(K_i\) acting on
four-vectors and verify
\[
 [J_i,J_j]=i\epsilon_{ijk}J_k,\quad
 [J_i,K_j]=i\epsilon_{ijk}K_k,\quad
 [K_i,K_j]=-i\epsilon_{ijk}J_k.
\]
Explain the minus sign in the last commutator.

\SupplementaryExercise{8}{Two Oscillator Families from One Complex Field}
{(Adapted from B. Mistlberger and K. Zhou, Stanford PHYS-330,
Problem Set 2, 2022, Problem 2(a), pp. 1--2.)}
{zhou-stanford-phys330-ps2-2022}
Quantize a free complex scalar field and show that its canonical algebra
requires two mutually commuting oscillator families \(a(\mathbf p)\) and
\(b(\mathbf p)\).  Identify the one-particle states they create.

\SupplementaryExercise{9}{Charge Counts Particles Minus Antiparticles}
{(Adapted from B. Mistlberger and K. Zhou, Stanford PHYS-330,
Problem Set 2, 2022, Problem 2(b)--(d), p. 2.)}
{zhou-stanford-phys330-ps2-2022}
Find the Noether charge for the phase symmetry of a complex scalar field.
Normal order it, determine its commutators with the field and the two
creation operators, and identify particle and antiparticle charges.

\SupplementaryExercise{10}{Energy Counts Both Charge Signs}
{(Adapted from B. Mistlberger and K. Zhou, Stanford PHYS-330,
Problem Set 2, 2022, Problem 2(e)--(f), p. 2.)}
{zhou-stanford-phys330-ps2-2022}
Write the normal-ordered Hamiltonian and momentum of a complex scalar field
in terms of \(a,b\).  Show that particles and antiparticles have opposite
charge but the same positive energy and four-momentum spectrum.

\SupplementaryExercise{11}{Causal Commutators of a Charged Scalar}
{(Adapted from B. Mistlberger and K. Zhou, Stanford PHYS-330,
Problem Set 2, 2022, Problem 2, pp. 1--2.)}
{zhou-stanford-phys330-ps2-2022}
For the complex scalar mode expansion, calculate
\([\Phi(x),\Phi(y)]\) and
\([\Phi(x),\Phi^\dagger(y)]\).  Show explicitly which one vanishes
identically and why the other vanishes at spacelike separation.

\SupplementaryExercise{12}{Boost Charges on Scalar Creation Operators}
{(Adapted from B. Mistlberger and K. Zhou, Stanford PHYS-330,
Problem Set 2, 2022, Problem 3, p. 3.)}
{zhou-stanford-phys330-ps2-2022}
Derive the normal-ordered scalar boost generator \(M^{0i}\) in momentum
space, including the term required by the momentum-space measure.  Check
that its commutator generates an infinitesimal Lorentz boost of a
relativistically normalized one-particle state.

\SupplementaryExercise{13}{Representation-Independent Gamma Traces}
{(Adapted from D. Tong, \emph{Quantum Field Theory: Example Sheet 3},
2006, Problem 3, pp. 1--2.)}{tong-qft-sheet3-2006}
Using only
\(\{\gamma^\mu,\gamma^\nu\}=2\eta^{\mu\nu}\), prove the traces of two
and four gamma matrices and show that every trace of an odd number of gamma
matrices vanishes.

\SupplementaryExercise{14}{Spinor Generators of Lorentz Transformations}
{(Adapted from D. Tong, \emph{Quantum Field Theory: Example Sheet 3},
2006, Problem 2, p. 1.)}{tong-qft-sheet3-2006}
Define
\(\mathcal J^{\mu\nu}=-\tfrac i4[\gamma^\mu,\gamma^\nu]\).
Verify that these matrices obey the Lorentz Lie algebra and prove that
\(D(\Lambda)^{-1}\gamma^\rho D(\Lambda)
=\Lambda^\rho{}_\sigma\gamma^\sigma\) infinitesimally.

\SupplementaryExercise{15}{Dirac Spinor Norms and Orthogonality}
{(Adapted from D. Tong, \emph{Quantum Field Theory: Example Sheet 3},
2006, Problem 4, p. 2.)}{tong-qft-sheet3-2006}
For on-shell positive- and negative-frequency Dirac spinors, derive their
ordinary and Dirac-adjoint norms and the orthogonality relations.  State
clearly which momentum reversal is needed for an ordinary Hermitian inner
product.

\SupplementaryExercise{16}{Completeness of Free Dirac Spinors}
{(Adapted from D. Tong, \emph{Quantum Field Theory: Example Sheet 3},
2006, Problem 5, pp. 2--3.)}{tong-qft-sheet3-2006}
Derive the positive- and negative-frequency spin sums in the conventions of
this volume and show that they project onto the two on-shell eigenspaces of
the Dirac operator.

\SupplementaryExercise{17}{Equal-Time Dirac Anticommutators}
{(Adapted from D. Tong, \emph{Quantum Field Theory: Example Sheet 3},
2006, Problem 6, p. 3.)}{tong-qft-sheet3-2006}
Use the free Dirac mode expansion and its spin sums to prove
\[
 \{\psi_\ell(\mathbf x,t),\psi_m^\dagger(\mathbf y,t)\}
 =\delta_{\ell m}\delta^3(\mathbf x-\mathbf y),
\]
with the other equal-time anticommutators zero.

\SupplementaryExercise{18}{What Fails with Bosonic Spinors}
{(Adapted from D. Tong, \emph{Quantum Field Theory: Example Sheet 3},
2006, Problem 8, p. 4.)}{tong-qft-sheet3-2006}
Attempt to impose bosonic commutators on the Dirac mode expansion while
retaining its equal-time field algebra.  Show that one oscillator family
must have the wrong sign and that the normal-ordered Hamiltonian is
unbounded below.

\SupplementaryExercise{19}{Proca Constraints and Degrees of Freedom}
{(Adapted from D. Tong, \emph{Lectures on Quantum Field Theory}, 2006,
Section 6.1.)}{tong-qft-free-fields-2006}
Starting from the Proca equation
\(\partial_\mu F^{\mu\nu}-m^2A^\nu=0\) in mostly-plus convention, derive
the transversality condition and the Klein--Gordon equation.  Count the
propagating polarizations.

\SupplementaryExercise{20}{Completeness of Massive Vector Polarizations}
{(Adapted from D. Tong, \emph{Lectures on Quantum Field Theory}, 2006,
Sections 6.1--6.2.)}{tong-qft-free-fields-2006}
Construct three polarization vectors for a massive vector with momentum
along the 3-axis and prove
\[
 \sum_{\sigma=-1}^{1}e^\mu(\mathbf p,\sigma)
 e^{\nu *}(\mathbf p,\sigma)
 =\eta^{\mu\nu}+\frac{p^\mu p^\nu}{m^2}.
\]

\SupplementaryExercise{21}{The Longitudinal Mode and a Conserved Current}
{(Adapted from D. Tong, \emph{Lectures on Quantum Field Theory}, 2006,
Sections 6.1--6.2.)}{tong-qft-free-fields-2006}
Find the high-energy behavior of the longitudinal massive-vector
polarization.  Show that its coupling to \(J_\mu\) has a smooth massless
limit when \(p_\mu J^\mu=0\), and determine the leading residual behavior.

\SupplementaryExercise{22}{Little-Group Translations Become Gauge Shifts}
{(Adapted from D. Tong, \emph{Lectures on Quantum Field Theory}, 2006,
Section 6.2.)}{tong-qft-free-fields-2006}
For a massless spin-one polarization, show that the translation part of the
little group changes \(e^\mu\) by a multiple of \(p^\mu\).  Prove that a
Lorentz-covariant emission amplitude is insensitive to this change exactly
when the current is conserved.

\SupplementaryExercise{23}{A Covariant Field Strength from Helicity States}
{(Adapted from D. Tong, \emph{Lectures on Quantum Field Theory}, 2006,
Sections 6.1--6.2.)}{tong-qft-free-fields-2006}
Define
\(f^{\mu\nu}(p,\sigma)=i(p^\mu e^\nu-p^\nu e^\mu)\)
for helicity \(\sigma=\pm1\).  Show that it is invariant under
\(e^\mu\to e^\mu+\alpha p^\mu\), obeys the source-free Maxwell identities,
and carries the two helicities in the
\((1,0)\oplus(0,1)\) representation.

\SupplementaryExercise{24}{The Transverse Canonical Commutator}
{(Adapted from D. Tong, \emph{Lectures on Quantum Field Theory}, 2006,
Section 6.2.)}{tong-qft-free-fields-2006}
In Coulomb gauge, derive
\[
 [A_i(\mathbf x,t),E_j(\mathbf y,t)]
 =i\left(\delta_{ij}-\frac{\partial_i\partial_j}{\nabla^2}\right)
 \delta^3(\mathbf x-\mathbf y).
\]
Explain how this projector implements the two physical photon degrees of
freedom.

\SupplementaryExercise{25}{Why the Annihilation Field Is Spatially Nonlocal}
{(Adapted from D. Harlow, \emph{Lecture Notes for Physics 8.323}, 2024,
Section 3.3, pp. 35--36.)}{harlow-qft1-fields-2024}
At fixed time, solve for the positive-frequency part \(\phi^{(+)}\) of a
real scalar field in terms of \(\phi\) and \(\pi\).  Show that the operator
\(\sqrt{m^2-\nabla^2}\) makes this separation nonlocal in space even though
the full field is local.

\SupplementaryExercise{26}{Antiparticles Restore Charged-Field Locality}
{(Adapted from D. Harlow, \emph{Lecture Notes for Physics 8.323}, 2024,
Section 3.5, pp. 37--38.)}{harlow-qft1-fields-2024}
Try to build a charged scalar field from particle annihilation operators and
their adjoints alone.  Show why a local field of definite charge instead
requires an independent antiparticle oscillator with the opposite charge.

\SupplementaryExercise{27}{Microcausality from the Pauli--Jordan Function}
{(Adapted from D. Harlow, \emph{Lecture Notes for Physics 8.323}, 2024,
Section 3.4, pp. 35--37.)}{harlow-qft1-fields-2024}
Write the free scalar commutator as the Pauli--Jordan integral.  Prove its
vanishing at spacelike separation by evaluating it at equal time in a
suitable Lorentz frame, and explain why the separate positive-frequency
function does not vanish.

\SupplementaryExercise{28}{Antiunitary CPT Reverses Operator Order}
{(Adapted from D. Harlow, \emph{Lecture Notes for Physics 8.323}, 2024,
Section 6.1, pp. 69--71.)}{harlow-qft1-fields-2024}
Let \(\Theta\) be antiunitary and implement CPT on local fields.  Show how a
vacuum correlator transforms, including complex conjugation and reversal of
operator order.  Explain why locality is needed to compare the transformed
correlator with the original one.

\SupplementaryExercise{29}{CPT of a Local Hermitian Interaction}
{(Adapted from D. Harlow, \emph{Lecture Notes for Physics 8.323}, 2024,
Section 6.1, pp. 69--71.)}{harlow-qft1-fields-2024}
Using the CPT rules for scalar, vector, and Dirac bilinears, prove that a
local Hermitian Lorentz-scalar interaction density satisfies
\(\Theta\mathcal H_I(x)\Theta^{-1}=\mathcal H_I(-x)\).
Deduce invariance of \(H_I(0)\).

\SupplementaryExercise{30}{The Locality Sign and Spin--Statistics}
{(Adapted from D. Harlow, \emph{Lecture Notes for Physics 8.323}, 2024,
Section 6.2, pp. 71--74.)}{harlow-qft1-fields-2024}
For a covariant free field of spin \(j\), compare its two-point function at
spacelike-separated points before and after a rotation that exchanges the
points.  Show how positivity and locality require the exchange sign
\((-1)^{2j}\), and state the resulting statistics assignment.
\end{enumerate}
